\documentclass[a4paper,10pt]{article}
\usepackage[margin=2.5cm]{geometry}
\usepackage[utf8]{inputenc}
\usepackage[colorlinks=true,urlcolor=blue]{hyperref}
\usepackage{amsmath}
\usepackage{graphicx}
\usepackage{float}
\usepackage{caption}
\usepackage{subcaption}

%\usepackage{listings} %Alternative to minted
\usepackage{minted}



\setlength{\parindent}{0em}
\setlength{\parskip}{1em}

\title{\textbf{1RT730} 
	\\ Report for Hand-in Assignment 2}
\author{Student Name}
\date{\today}

\begin{document}
	
	\maketitle
	
	\section{Introduction}
    Motivation and relevance of trustworthiness in LLMs. Mention about the benchmark selected.

    \section{Model and set-up}
    Model description, dataset subsets and the three trust dimensions.

    \section{Results - Part A (Mandatory)}
    Benchmark dataset evaluation results, with metrics plots, tables, qualitative examples, etc. 

    \section{Results - Part B (Bonus)}
    Motivation of selected mitigation methods, implementation details, post-mitigation evaluation and comparison.  

    \section{Discussion and conclusions}
    Ethical implications, observed limitations, and suggestions for future work, if any. Conclude your findings

	

\hfill \break
\textit{\textbf{Use of generative AI}}

Write a few sentences if you have or have not used any generative AI, and if so how.

\begin{thebibliography}{1}

% Include references for all tools, datasets, and key literature used.

	%\bibitem{Item} Add here.
		
\end{thebibliography}


\appendix
\section{Code}
Here you can insert your code. Or you could include all your codes in a link (GitHub/HuggingFace repo or similar platform).


Please refer to the Assignment instructions on Studium. 


\label{app:code}
\inputminted[frame=lines,framesep=2mm,linenos]{python}{your_code.py}

	
\end{document}